\documentstyle[%


righttag%


,11pt%


,amssymb%


,russian%               remove in Eng.


,izvru%                 Set izven% in Eng.


% ,izvru-ks             for brief communications


]{amsart}





% {,} {\leq} {<}





\newcommand{\tg}{\operatorname{tg}}


\newcommand{\tts}[1]{}





%\hoffset=-15mm%                  For Eng.


\setcounter{page}{3}


\god{2005}


\nomer{2~(513)} %                        Remove in Eng.


%\nomer{Vol.\,49, No.\,2}%               Set in Eng.


%\str{\pageref{begin}--\pageref{end}}%  Set in Eng.


\udk{517.962}





\author{\it А.Ю.\,Александров, А.П.\,Жабко}


% NB:   ^^ remove in Eng.


\title{Об устойчивости решений\\ нелинейных разностных систем}





\date{}


%\pagestyle{myheadings}%         Set in Eng.


\pagestyle{plain} %              Remove in Eng.


\begin{document}


%\thispagestyle{plain}%          Set in Eng.


\maketitle


\label{begin}











Уравнения в конечных разностях широко применяются при описании


динамических систем, состояния которых известны (измеряются) в дискретные


моменты времени. К таким системам относятся, например, системы управления


с дискретными регуляторами [1]. Разностные уравнения являются основным


математическим аппаратом при изучении нелинейных импульсных систем [2].


Численное решение уравнений различных типов также приводит к замене


непрерывных систем дискретными [1].





Одно из направлений исследований, возникающих в указанных приложениях


разностных уравнений, связано с анализом устойчивости их решений.


В частности, задачи сходимости итерационных процессов --- это 


фактически задачи


устойчивости разностных систем. Методы исследования устойчивости


хорошо разработаны для линейных разностных уравнений.


Для нелинейных систем, как и для систем обыкновенных дифференциальных


уравнений, доказана теорема об устойчивости по линейному приближению


([2], c.\,38--40). Наиболее общим методом анализа устойчивости решений


разностных систем является дискретный аналог второго метода Ляпунова


([2], c.\,20--35).


В данной работе он используется для получения условий


асимптотической устойчивости по нелинейному приближению. При этом 


в качестве системы первого приближения рассматриваются уравнения 


с однородными правыми частями.


\medskip





{\bf 1. Предварительное обсуждение.} В данном разделе


проводится сравнение однородной разностной системы и соответствующей ей


системы обыкновенных дифференциальных уравнений в смысле устойчивости по


Ляпунову.





Рассмотрим разностную систему


\begin{equation}


\bold X(k+1)=\bold X(k)+\bold F(\bold X(k))


\end{equation}


и однородную систему дифференциальных уравнений


\begin{equation}


\Dot{\bold Z} =\bold F(\bold Z),


\end{equation}


где $\bold X(k)$, $\bold Z$ --- $n$-мерные векторы, элементы векторной


функции $\bold F(\bold Z)$ определены и непрерывно дифференцируемы


при всех $\bold Z\in \bold E^n$ и являются однородными


функциями порядка $\mu>1$.


Будем считать, что целочисленный аргумент $k$ во всех 


разностных уравнениях работы принимает значения $0,1,\dotsc$


Из однородности $\bold F(\bold Z)$


следует, что системы (1) и (2) имеют нулевые решения.





\begin{lem}


Если нулевое решение системы $(2)$ асимптотически устойчиво,


то нулевое решение системы $(1)$ также асимптотически устойчиво.


\end{lem}





\begin{pf}


Известно ([3], c.\,115--123), что для


системы (2), имеющей асимптотически устойчивое нулевое решение,


существуют


функции Ляпунова $V(\bold Z)$ и $W(\bold Z)$, которые обладают следующими


свойствами:


\begin{itemize}


\item[1)] $V(\bold Z)$, $W(\bold Z)$ --- непрерывно дифференцируемые


однородные функции порядка


$m$ и\linebreak $m+\mu-1$ соответственно, $m>1$;


\item[2)] $V(\bold Z)$, $W(\bold Z)$ --- положительно-определенные


функции;


\item[3)] функция $V(\bold Z)$ удовлетворяет уравнению


\end{itemize}


$$


\bigg(\frac {\partial V(\bold Z)}{\partial \bold Z}


\bigg)^\top \bold F(\bold Z)


= -W(\bold Z).


$$





Рассмотрим приращение функции $V(\bold Z)$


на решениях системы (1). Получим


$$


\Delta V=V(\bold X(k+1))-V(\bold X(k))


=-W(\bold X(k))+W_1(\bold X(k)).


$$


Здесь $W_1(\bold Z) / \|\bold Z\|^{m+\mu-1}\to 0$ при $\|\bold Z\|\to 0$


($\|\cdot\|$ --- евклидова норма вектора).


Следовательно ([3], c.\,118), существуют такие положительные постоянные


$b_1$, $b_2$ и $\delta$, что при $\| \bold X(k)\|< \delta$


справедливы оценки


\begin{equation}


-b_1 \|\bold X(k)\|^{m+\mu-1}\leq \Delta V\leq


-b_2 \|\bold X(k)\|^{m+\mu-1}.


\end{equation}


Значит ([2], c.\,27--30), нулевое решение системы (1) асимптотически


устойчиво.


\end{pf}





Обратное утверждение, как показывает следующий пример, неверно.





\begin{exam}


Рассмотрим разностную систему


\begin{equation}


\begin{split}


x_1(k+1)&= x_1(k)+ \rho^2_k \big(\alpha_k x_1(k)-


\beta_k x_2(k) \big),\\


x_2(k+1)&= x_2(k)+ \rho^2_k \big( \beta_k x_1(k)+


\alpha_k x_2(k) \big).


\end{split}


\end{equation}


Здесь $\rho$ и $\varphi$ --- полярные координаты


($x_1= \rho \cos \varphi$, $x_2= \rho \sin \varphi$),


$\rho_k=\rho(k)$, $\varphi_k=\varphi(k)$, $\alpha_k=\alpha(\varphi_k)$,


$\beta_k=\beta(\varphi_k)$,


а функции $\alpha(\varphi)$ и $\beta(\varphi)$ определяются по формулам


$$


\alpha(\varphi)=\frac{\beta(\varphi) \cos(10 \varphi)}{1+0{,}1


\sin(10 \varphi)},\quad


\beta(\varphi)=2-\sin(10 \varphi).


$$ 





Система (4) является однородной порядка $\mu=3$, причем ее правые части


дважды непрерывно дифференцируемы. Покажем, что нулевое решение этой 


системы асимптотически устойчиво.





В полярных координатах уравнения (4) принимают вид


\begin{equation}


\begin{array}{c}


\rho_{k+1}=\rho_{k} \sqrt{{(1+\alpha_k \rho^2_k)}^2+


{(\beta_k \rho^2_k)}^2}, \\[3mm]


\tg (\varphi_{k+1}-\varphi_k)= 


\dfrac{\beta_k \rho^2_k}{1+\alpha_k \rho^2_k}.


\end{array}


\end{equation}





Введем вспомогательную функцию


$$


c(\varphi)=\int_0^\varphi \frac{\alpha(\tau)}{\beta(\tau)}\, d\tau=


\ln(1+0{,}1 \sin(10 \varphi)).


$$





Пусть


\begin{multline*}


V(\rho, \varphi)=\rho^2 \exp(-2\, c(\varphi))


\bigg(1-\rho^2 \exp(-2\, c(\varphi)) \bigg(\gamma \varphi+ \\


+\int_0^\varphi \exp(2 c(\tau)) \beta(\tau) (1+(c'(\tau))^2-


c''(\tau)) d\tau\bigg)\bigg),


\end{multline*}


где


$$


\gamma =-\frac1{2\pi} \int_0^{2\pi} \exp(2 c(\tau)) \beta(\tau)


(1+(c'(\tau))^2-c''(\tau)) d\tau >0.


$$





Приращение функции $V(\rho, \varphi)$ на решениях системы (5)


может быть представлено в виде


$$


\Delta V=V(\rho_{k+1}, \varphi_{k+1})- V(\rho_k, \varphi_k)=


-\gamma \rho^6_k \beta_k 


\exp(-4\, c(\varphi_k)) + W_1(\rho_k, \varphi_k).


$$


Здесь функция $W_1(\rho, \varphi)$ в достаточно малой окрестности точки


$\rho=0$ и при всех $\varphi\in (-\infty, +\infty)$


удовлетворяет неравенству $|W_1(\rho, \varphi)|\leq b_1 \rho^8$,


где $b_1$ --- положительная постоянная.





Значит, существуют числа


$\delta>0$ и $b_2>0$ такие, что при $\rho_k< \delta$,


$\varphi_k\in (-\infty, +\infty)$ справедлива оценка


$\Delta V\leq -b_2 \, \rho^6_k$,


из которой и следует асимптотическая


устойчивость нулевого решения системы~(4).





Соответствующая уравнениям (4) система


обыкновенных дифференциальных уравнений в полярных координатах имеет вид


\begin{equation}


\Dot \rho=\alpha(\varphi) \rho^3, \quad


\Dot \varphi=\beta(\varphi) \rho^2.


\end{equation}


При этом $\int\limits_0^{2\pi} \alpha(\varphi)/\beta(\varphi) d\varphi=0$.


Значит, нулевое решение системы (6) устойчиво, но не является 


асимптотически устойчивым.


\end{exam}





Следующий пример показывает, что если


нулевое решение разностной системы вида~(1) неустойчиво, то


соответствующая ей система~(2) может иметь устойчивое нулевое решение.





\begin{exam}


Рассмотрим систему


\begin{equation}


\begin{split}


x_1(k+1)&= x_1(k)-\big(x^2_1(k)+x^2_2(k)\big)x_2(k),\\


x_2(k+1)&= x_2(k)+\big(x^2_1(k)+x^2_2(k)\big)x_1(k).


\end{split}


\end{equation}


В качестве функции Ляпунова выберем функцию


$V(x_1, x_2)=x^2_1+x^2_2$. Получим


$$


V(x_1(k+1), x_2(k+1))-V(x_1(k), x_2(k))={(x^2_1(k)+x^2_2(k))}^3.


$$


Отсюда вытекает ([4], c.\,82), что нулевое решение уравнений~(7)


неустойчиво.





В то же время нулевое решение соответствующей однородной системы


обыкновенных дифференциальных уравнений


$$


\Dot z_1=-(z^2_1+z^2_2)z_2, \quad \Dot z_2=(z^2_1+z^2_2)z_1


$$


является устойчивым.


\end{exam}








{\bf 2. Оценки решений нелинейных разностных систем.}


Известно ([3], с.\,117), что если нулевое решение системы (2)


асимптотически устойчиво, то при всех


$\bold Z_0\in \bold E^n$, $t_0\in (-\infty, +\infty)$, $t\geq t_0$


справедливы неравенства


$$


c_1 \|\bold Z_0\|


\big(1 +c_2 \|\bold Z_0\|^{\mu-1}(t-t_0)\big)^{-\frac1{\mu-1}}


\leq \| \bold Z(t,\bold Z_0,t_0)\|\leq 


%%


c_3 \|\bold Z_0\|


\big(1 +c_4 \|\bold Z_0\|^{\mu-1}(t-t_0)\big)^{-\frac1{\mu-1}},


$$


где $c_1$, $c_2$, $c_3$, $c_4$ --- положительные постоянные,


$\bold Z(t,\bold Z_0,t_0)$ --- решение уравнений (2), проходящее при


$t=t_0$ через точку $\bold Z_0$.





Оценим скорость стремления к началу координат решений разностной


системы (1) в случае, когда для нее построена однородная функция


Ляпунова, удовлетворяющая условиям теоремы об асимптотической


устойчивости ([2], c.\,27--30).





Пусть нулевое решение системы дифференциальных уравнений (2)


асимптотически устойчиво. Из доказательства леммы~1 следует, что


тогда для любого решения $\bold X(k)$ системы (1), начинающегося при


$k=k_0\geq 0$ в достаточно малой окрестности точки $\bold X=\bold 0$,


при всех $k\geq k_0$ выполняются неравенства~(3).





Известно ([3], c.~118), что если $V(\bold X)$ ---


положительно-определенная однородная порядка $m$ функция, то для любых


$\bold X\in \bold E^n$ справедливы оценки


\begin{equation}


a_1 \|\bold X\|^{m}\leq V(\bold X)\leq a_2 \|\bold X\|^{m},


\end{equation}


где $a_1$, $a_2$ --- положительные постоянные.





Используя неравенства (8), получаем, что при $k=k_0, k_0+1,\dotsc$


имеют место соотношения


\begin{equation}


V (\bold X(k))-b_1 \bigg(\frac{V (\bold X(k))}{a_1}\bigg)^{\frac{\mu-1}m+1}


\leq V(\bold X(k+1))\leq 


V (\bold X(k))-b_2


\bigg(\frac{V (\bold X(k))}{a_2}\bigg)^{\frac{\mu-1}m+1}.


\end{equation}








\begin{lem}


Пусть для членов последовательности $\{v_k\}$ выполнены неравенства


\begin{equation}


0\leq v_{k+1}\leq v_k-av^\lambda_k, \quad k=0,1,\dotsc,


\end{equation}


где $a>0$, $\lambda>1$, $v_0\geq 0$, $a \lambda v^{\lambda-1}_0\leq 1$.


Тогда при всех $k=0,1,\dotsc$ справедлива оценка


\begin{equation}


v_k\leq v_0 \big(1+ a (\lambda-1) v^{\lambda-1}_0 k


\big)^{-\frac1{\lambda-1}}.


\end{equation}


\end{lem}





\begin{pf}


Применяем метод математической индукции.


При $k=0$ неравенство (11) верно. Предположим, что оно имеет место для всех


$k\leq l$. Тогда $v_l\leq v_0$, откуда следует


$a \lambda v^{\lambda-1}_l{\leq}1$.


Функция $v_l-av^\lambda_l$ монотонно возрастает на промежутке


$[0, (a\lambda)^{-1/(\lambda-1)}]$. Значит, справедливы соотношения


$$


v_{l+1}\leq v_l (1-av^{\lambda-1}_l)\leq


v_0 (1+ \theta l)^{-\frac1{\lambda-1}}


\bigg(1-\frac{\theta}{(\lambda-1)(1+ \theta l)}\bigg)=


v_0 \big(1+ \theta (l+1)\big)^{-\frac1{\lambda-1}} g(\theta),


$$


где


$$


\theta=a (\lambda-1) v^{\lambda-1}_0, \quad


g(\theta)=\bigg(1+\frac{\theta}{1+ \theta l}\bigg)^{\frac1{\lambda-1}}


\bigg(1-\frac{\theta}{(\lambda-1)(1+ \theta l)}\bigg).


$$





Заметим теперь, что $g(0)=1$ и при этом


функция $g(\theta)$ монотонно убывает на промежутке


$[0, +\infty)$. Следовательно, неравенство (11) выполняется и при


$k=l+1$.


\end{pf}





Аналогичным образом доказывается





\begin{lem}


Пусть для членов последовательности $\{v_k\}$ выполнены неравенства


\begin{equation}


v_k-bv^\lambda_k\leq v_{k+1}\leq


(b\lambda)^{-\frac1{\lambda-1}}, \quad k=0,1,\dotsc,


\end{equation}


где $b>0$, $\lambda>1$, $v_0\geq 0$, $2b \lambda v^{\lambda-1}_0\leq 1$.


Тогда при всех $k=0,1,\dotsc$ имеет место оценка


\begin{equation}


v_0 \big(1+ 2b (\lambda-1) v^{\lambda-1}_0 k


\big)^{-\frac1{\lambda-1}}\leq v_k.


\end{equation}


\end{lem}





\begin{note}


Если неравенства (10) (неравенства (12))


выполняются лишь при $k=0,1,...,N$,     %%%%%% <-- restore \dotsc


то оценка (11) (оценка (13))


будет справедлива при $k=0,1,\dotsc,N+1$.


\end{note}





С помощью лемм 2, 3 получаем, что из


неравенств (8), (9) следует существование


числа $\delta>0$ такого, что для любого решения


$\bold X(k)$ системы (1) с начальными данными $\bold X(k_0)=\bold X_0$,


удовлетворяющими условиям $k_0\geq 0$, $\|\bold X_0\|<\delta$,


при всех $k\geq k_0$ имеют место оценки


\begin{equation}


c_1 \|\bold X_0\|


\big(1 +c_2 \|\bold X_0\|^{\mu-1}(k-k_0)\big)^{-\frac1{\mu-1}}


\leq \| \bold X(k)\|\leq 


c_3 \|\bold X_0\|


\big(1 +c_4 \|\bold X_0\|^{\mu-1}(k-k_0)\big)^{-\frac1{\mu-1}}.


\end{equation}


Здесь $c_1=(a_1/a_2)^{1/m}$, $c_2=2(\mu-1)b_1/(ma_1)$, $c_3=1/c_1$,


$c_4=(\mu-1)b_2/(ma_2)$.


\medskip





{\bf 3. Анализ устойчивости разностных систем по нелинейному


приближению.} Далее будем предполагать, что для системы (1) существуют


однородные функции Ляпунова $V(\bold X)$ и $W(\bold X)$,


которые обладают свойствами, указанными в доказательстве леммы~1.


Следовательно, нулевое решение этой системы асимптотически устойчиво.





Наряду с уравнениями (1) рассмотрим возмущенные уравнения


\begin{equation}


\bold X(k+1)=\bold X(k)+\bold F(\bold X(k))+\bold R(k,\bold X(k)).


\end{equation}


Здесь векторная функция $\bold R(k,\bold X)$ определена при


$k=0,1,\dotsc$ и $\|\bold X\|\leq h$ ($h$ --- положительная


постоянная), непрерывна по $\bold X$ и удовлетворяет условию


$\|\bold R(k,\bold X)\|\leq c \|\bold X\|^\sigma$, где


$c>0$, $\sigma>0$.





\begin{thm}


При выполнении неравенства $\sigma>\mu$


нулевое решение системы $(15)$ асимптотически устойчиво.


\end{thm}





\begin{pf}


Приращение функции $V(\bold X)$ на решениях


возмущенной системы можно представить в виде


\begin{multline}


\Delta V=V\big(\bold X(k)+\bold F(\bold X(k))\big)


-V(\bold X(k))+ \\


+\bigg(\frac {\partial V}{\partial \bold X}


\big(\bold X(k)+\bold F(\bold X(k))+\theta_k \bold R(k,\bold X(k))\big)


\bigg)^\top \bold R(k,\bold X(k)),


$$


\end{multline}


где $\theta_k\in (0, 1)$. Значит, при $\|\bold X(k)\|\leq h$ справедлива


оценка


$$


\Delta V\leq -b_1\, \|\bold X(k)\|^{m+\mu-1}+\psi(\bold X(k))


\|\bold X(k)\|^{m+\mu-1} +


b_2 \|\bold X(k)\|^{\sigma}(


\|\bold X(k)\|+\|\bold X(k)\|^{\mu}+\|\bold X(k)\|^{\sigma})^{m-1}.


$$


Здесь $b_1$, $b_2$ --- положительные постоянные, а функция


$\psi(\bold X)$ неотрицательна и стремится к нулю при $\|\bold X\|\to 0$.





Если $\sigma>\mu$, то для достаточно малых значений


$\|\bold X(k)\|$ и при всех $k=0,1,\dotsc$ имеет место неравенство


$\Delta V\leq -b_1 \|\bold X(k)\|^{m+\mu-1}/2$.


Следовательно, функция $V(\bold X)$ удовлетворяет требованиям теоремы об


асимптотической устойчивости ([2], c.\,27--30).


\end{pf}





\begin{note}


Аналогичным образом можно показать, что если


$\|\bold X(k)\|$ достаточно мала, то при всех $k=0,1,\dotsc$ выполнено


неравенство $\Delta V\geq -b_3 \|\bold X(k)\|^{m+\mu-1}$, где $b_3>0$. 


Значит, для решений системы (15) в некоторой окрестности


точки $\bold X=\bold 0$ справедливы те же самые оценки (14), что и для


решений невозмущенной системы, но, вообще говоря, с другими значениями


постоянных $c_1$, $c_2$, $c_3$, $c_4$.


\end{note}





\begin{note}


Предположим, что условие 2), %%%% condition?? cf. proof of lemma 1


которому удовлетворяют функции $V(\bold X)$ и


$W(\bold X)$, заменено условием


\begin{itemize}


\item[2$'$)] функция $V(\bold X)$ не является знакопостоянной отрицательной,


а функция $W(\bold X)$ отрицательно определена.


\end{itemize}


Тогда, применяя первую теорему о неустойчивости ([4], c.\,82),


можно показать, что нулевое решение системы (1) неустойчиво, а при


$\sigma>\mu$ неустойчивость нулевого решения сохраняется и для


возмущенной системы.


\end{note}





{\bf 4. Исследование устойчивости систем, находящихся под воздействием


неограниченных возмущений.} Пусть теперь векторная функция


$\bold R(k,\bold X)$ при всех $k=0,1,\dotsc$ и $\|\bold X\|\leq h$


удовлетворяет условию


\begin{equation}


\|\bold R(k,\bold X)\|\leq c (k+1)^\alpha \|\bold X\|^\sigma,


\end{equation}


где $c$, $\alpha$, $\sigma$ --- положительные постоянные.


Таким образом, возмущения в системе (15), вообще говоря, 


являются неограниченными, а параметр $\alpha$ определяет степень их


роста.





Исследуем вопрос, насколько порядок возмущений (величина~$\sigma$)


должен превышать порядок однородности правых частей системы (1),


чтобы из асимптотической устойчивости


нулевого решения этой системы следовала асимптотическая устойчивость


нулевого решения возмущенных уравнений.





\begin{thm}


При выполнении неравенства


$\sigma>\mu+\alpha(\mu-1)$


нулевое решение системы~$(15)$ асимптотически устойчиво.


\end{thm}





\begin{pf}


Используя представление приращения функции $V(\bold X)$ на 


решениях возмущенных уравнений в виде (16) и условие (17),


получаем, что при $\|\bold X(k)\|<\delta$, где $\delta$ --- 


достаточно малое


положительное число, и для всех $k=0,1,\dotsc$ справедлива оценка


$$


\Delta V\leq -b_1 \|\bold X(k)\|^{m+\mu-1}+


b_2 (k+1)^\alpha \|\bold X(k)\|^{m+\sigma-1}+


b_3 (k+1)^{\alpha m} \|\bold X(k)\|^{\sigma m},


$$


где $b_1, b_2, b_3>0$.





Покажем, что существуют числа $\gamma$, $A$ и $L$


такие, что для любого решения $\bold X(k)$


системы (15) с начальными данными $\bold X(k_0)=\bold X_0$,


удовлетворяющими условиям


\begin{equation}


k_0\geq L,\quad \| \bold X_0\|^{\mu-1}<{\gamma}/{k_0},


\end{equation}


при всех $k\geq k_0$ имеет место неравенство


\begin{equation}


\| \bold X(k)\|^{\mu-1}<{A}/k.


\end{equation}





Пусть $a_1=\min\limits_{\|\bold X\|=1} V(\bold X)$,


$a_2=\max\limits_{\|\bold X\|=1} V(\bold X)$.


Выберем числа $\gamma$ и $A$ так, чтобы выполнялись соотношения


$$


0<\gamma<\frac{2m a_2}{(\mu-1)b_1}, \quad A> \frac{2m a_2}{(\mu-1)b_1}


\bigg(\frac{a_2}{a_1}\bigg)^{\frac{\mu-1}{m}}.


$$


Далее задаем $L>0$ настолько большим, чтобы при


$k\geq L$ были справедливы неравенства


$$


{A}/k<\delta^{\mu-1},\ \ 


k\geq \frac{(\mu-1+m) b_1 \gamma}{2 m a_2},\ \ 


4 b_2 (k+1)^{\alpha} ({A}/{k})^{\frac{\sigma-\mu}{\mu-1}}


<b_1,\ \ 4 b_3 (k+1)^{\alpha m}


({A}/{k})^{\frac{(\sigma-1)m}{\mu-1}-1} <b_1.


$$





Рассмотрим решение $\bold X(k)$ системы (15) с начальными данными,


удовлетворяющими условиям (18). Пусть существует такое натуральное число


$k_1>k_0$, что


$\left\| \bold X(k_1)\right\|^{\mu-1}\geq A/k_1$,


а при $k=k_0,\dotsc, k_1-1$ имеет место оценка~(19).


Тогда для всех $k=k_0,\dotsc, k_1-1$ получаем


$$


\Delta V\leq -\frac{b_1}2\|\bold X(k)\|^{m+\mu-1}\leq


-\frac{b_1}{2}\bigg(\frac{V(\bold X(k))}{a_2}\bigg)^{\frac{\mu-1}m+1}.


$$


Заметим, что при этом


$$


\frac{(\mu-1+m)b_1}{2 m a_2}


\bigg(\frac{V (\bold X_0)}{a_2}\bigg)^{\frac{\mu-1}m}<


\frac{(\mu-1+m)b_1 \gamma}{2 m a_2 k_0}\leq 1.


$$





Применяя лемму 2, приходим к неравенству


$$


V (\bold X(k_1))\leq V (\bold X_0)\bigg(1+


\frac{(\mu-1) b_1}{2m} a^{-\frac{\mu-1}m-1}_2


V^{\frac{\mu-1}m} (\bold X_0) (k_1-k_0)


\bigg)^{-\frac{m}{\mu-1}}.


$$


Учитывая оценки (8), которым удовлетворяет функция $V (\bold X)$,


получаем неравенство


$$


k_1\bigg(\frac{(\mu-1)b_1}{2ma_2}-\frac{1}{A}\bigg(\frac{a_2}{a_1}


\bigg)^{\frac{\mu-1}{m}}\bigg)\leq k_0\bigg(\frac{(\mu-1)b_1}{2ma_2}


-\frac{1}{\gamma}\bigg).


$$


Из условий выбора чисел $\gamma$ и $A$


следует, что левая часть данного неравенства положительна,


а правая --- отрицательна. Приходим к противоречию. Значит, для решения


$\bold X(k)$ при всех $k\geq k_0$ имеет место оценка~(19).





Используя доказанное свойство решений системы (15),


а также их непрерывную зависимость от начальных данных,


получаем утверждение теоремы.


\end{pf}





{\bf 5. Построение нестационарных функций Ляпунова для нелинейных


систем.} Теорема~1 утверждает, что возмущения не нарушают 


асимптотической устойчивости нулевого решения, если их порядок 


выше порядка однородности функций, входящих в правые части


уравнений~(1). Покажем теперь, что для некоторых классов 


нестационарных возмущений асимптотическая устойчивость сохраняется 


и в случае, когда $\sigma\leq \mu$.





Пусть система (15) имеет вид


\begin{equation}


\bold X(k+1)=\bold X(k)+\bold F(\bold X(k))+\bold B_k 


\bold Q(\bold X(k)).


\end{equation}


Здесь $\bold B_k$ --- постоянные матрицы размерности $n\times l$,


а элементы $l$-мерного вектора $\bold Q(\bold X)$ представляют 


собой непрерывно дифференцируемые однородные функции порядка


$\sigma>1$.





Предположим, что последовательность $\bold B_k$ является ограниченной.


Кроме того, в дополнение к условиям 1)--3) на функции $V(\bold X)$ и


$W(\bold X)$, будем считать, что $V(\bold X)$ --- дважды непрерывно


дифференцируемая функция. Заметим, что для существования таких функций


Ляпунова достаточно, чтобы нулевое решение системы (2) было 


асимптотически устойчивым,


а векторная функция $\bold F(\bold X)$ дважды непрерывно


дифференцируемой ([3], c.\,119--123).





Рассмотрим последовательность $n\times l$-матриц


\begin{equation}


\bold I_0=\bold 0, \quad \bold I_k=\sum_{j=0}^{k-1} \bold B_j,


\quad k=1,2,\dotsc


\end{equation}





\begin{thm}


Если последовательность $(21)$ ограничена,


то при выполнении неравенства


\begin{equation}


2 \sigma>\mu+1


\end{equation}


нулевое решение системы $(20)$ асимптотически устойчиво.


\end{thm}





\begin{pf}


Функцию Ляпунова выбираем в виде


\begin{equation}


V_1(k,\bold X)=V(\bold X)-


\bigg(\frac{\partial V(\bold X)}{\partial \bold X}\bigg)^\top \bold I_k


\bold Q(\bold X).


\end{equation}


Получим


\begin{multline*}


\Delta V_1=V_1(k+1, \bold X(k+1))-V_1(k, \bold X(k))=-W(\bold X(k))+\\


%


+\frac12 \big(\bold X(k+1)-\bold X(k)


\big)^\top \frac {\partial^2 V(\bold Z_k)}


{\partial \bold X^2}\big(\bold X(k+1)-\bold X(k) \big)+\\


%


+\bigg(\frac{\partial V(\bold X(k))}{\partial \bold X}\bigg)^\top


\bold I_{k+1} \bold Q(\bold X(k))-\bigg(\frac{\partial


V(\bold X(k+1))}{\partial \bold X}\bigg)^\top


\bold I_{k+1} \bold Q(\bold X(k+1)),


\end{multline*}


где $\bold Z_k= \bold X(k)+\theta_k (\bold X(k+1)-\bold X(k))$,


$\theta_k\in (0, 1)$. Следовательно, при всех $k=0,1,\dotsc$ и


$\bold X(k)\in \bold E^n$ справедливы соотношения


\begin{gather*}


a_1 \|\bold X(k)\|^{m}-a_3


\|\bold X(k)\|^{m+\sigma-1}\leq V_1(k, \bold X(k))\leq


a_2 \|\bold X(k)\|^{m}+a_3\|\bold X(k)\|^{m+\sigma-1},\\


%


\Delta V_1\leq -b_1 \|\bold X(k)\|^{m+\mu-1}+


b_2 (\|\bold X(k)\|^{2 \mu}+ \|\bold X(k)\|^{2 \sigma})


(\|\bold X(k)\|+ \\


%


+\|\bold X(k)\|^{\mu}+\|\bold X(k)\|^{\sigma})^{m-2}


+b_3 (\|\bold X(k)\|^{\mu}+ \|\bold X(k)\|^{ \sigma})


(\|\bold X(k)\|+ \|\bold X(k)\|^{\mu}+\|\bold X(k)\|^{\sigma}


)^{m+\sigma-2},


\end{gather*}


где $a_1$, $a_2$, $a_3$, $b_1$, $b_2$, $b_3$ ---


положительные постоянные.





Если выполнено условие (22), то для достаточно малых значений


$\|\bold X(k)\|$ и при всех $k=0,1,\dotsc$ имеют место неравенства


$$


\frac{a_1}2 \|\bold X(k)\|^{m}


\leq V_1(k, \bold X(k))\leq 2 a_2 \|\bold X(k)\|^{m}, \quad


\Delta V_1\leq -\frac{b_1}2 \|\bold X(k)\|^{m+\mu-1}.


$$


Таким образом, функция $V_1(k,\bold X)$ удовлетворяет требованиям 


теоремы об асимптотической устойчивости ([2], c.\,27--30).


\end{pf}





Предположим теперь, что последовательность (21)


не является ограниченной. Однако существует такое число 


$0< \beta \leq 1$, что


\begin{equation}


\lim_{k\to \infty} \frac1{k^\beta} \bold I_k=\bold 0.


\end{equation}





\begin{thm}


При выполнении неравенства


$2\sigma\geq\mu+1+\beta(\mu-1)$


нулевое решение системы~$(20)$ асимптотически устойчиво.


\end{thm}





\begin{pf}


В качестве функции Ляпунова снова выберем


функцию (23). Для достаточно малых значений


$\|\bold X(k)\|$ и при всех $k=0,1,\dotsc$ получим оценки


\begin{gather*}


a_1 \|\bold X(k)\|^{m}-\omega_k (k+1)^\beta


\|\bold X(k)\|^{m+\sigma-1}\leq V_1(k, \bold X(k))\leq\\


%


\leq a_2 \|\bold X(k)\|^{m}+\omega_k (k+1)^\beta


\|\bold X(k)\|^{m+\sigma-1},\\


%


\Delta V_1\leq -b \|\bold X(k)\|^{m+\mu-1}+ \varepsilon_k


(k+1)^\beta (\|\bold X(k)\|^{m+ \mu+\sigma-2}+


\|\bold X(k)\|^{m+2 \sigma-2}),


\end{gather*}


где $a_1, a_2, b>0$, а последовательности неотрицательных чисел


$\omega_k$, $\varepsilon_k$ стремятся к нулю при $k\to \infty$.





Далее аналогично доказательству теоремы~2 можно доказать


существование положительных постоянных $\gamma$, $A$ и $L$ таких, что если


начальные данные решения $\bold X(k)$


системы~(20) удовлетворяют условиям (18), то при всех $k\geq k_0$


выполнено неравенство~(19).


\end{pf}








\begin{note}


При $0< \beta \leq 1$ имеем $(\mu+1+\beta(\mu-1))/2\leq \mu$. 


Следовательно, для всех рассматриваемых значений


параметра~$\beta$ теорема~4 уточняет полученный в разделе~3 критерий


асимптотической устойчивости по


нелинейному приближению: $\sigma>\mu$.


\end{note}





\begin{note}


Если последовательность (21) ограничена, то


предельное соотношение (24) справедливо


для любого $\beta>0$. Значит, для возмущений такого рода теорема~4 


определяет то же самое условие асимптотической устойчивости,


что и теорема~3, т.\,е. $2\sigma>\mu+1$.


\end{note}





\begin{exam}


Рассмотрим скалярное уравнение


\begin{equation}


x(k+1)=x(k)-a\, x^9(k)+b_k x^\sigma(k),


\end{equation}


где $a$ --- положительная постоянная, $\sigma$ --- 


рациональное число с нечетным знаменателем, $\sigma>1$.





Пусть $b_k=\sin k$, $k=0,1,\dotsc$ Тогда последовательность $\{I_k\}$, 


построенная по формулам (21), будет ограниченной. Применяя теорему~3, 


находим достаточное условие


асимптотической устойчивости нулевого решения уравнения~(25): $\sigma>5$.





Предположим теперь, что $b_k=\cos\pi \sqrt{k}$, $k=0,1,\dotsc$


В этом случае последовательность $\{I_k\}$ уже не является ограниченной, 


а предельное соотношение (24) справедливо тогда и только тогда, 


когда $\beta>1/2$. В соответствии с теоремой~4 для 


асимптотической устойчивости нулевого решения рассматриваемого уравнения 


достаточно выполнения неравенства $\sigma>7$.


\end{exam}








{\bf 6. Асимптотическая устойчивость по отношению к части переменных.}


Рассмотрим систему


\begin{equation}


\begin{split}


\bold X(k+1)&=\bold X(k)+\bold F(\bold X(k))+


\bold R(k,\bold X(k),\bold Y(k)), \\


\bold Y(k+1)&=\bold G(k,\bold Y(k))+


\bold D(k,\bold X(k),\bold Y(k)).


\end{split}


\end{equation}


Здесь $\bold X(k)$ и $\bold Y(k)$ --- векторы размерности $n$ и $l$


соответственно; элементы векторной функции $\bold F(\bold X)$ 


определены и непрерывно дифференцируемы


при всех $\bold X\in \bold E^n$ и являются однородными


функциями порядка $\mu>1$;


векторные функции $\bold G(k,\bold Y)$, $\bold R(k,\bold X,\bold Y)$,


$\bold D(k,\bold X,\bold Y)$ заданы при


$k=0,1,\dotsc$, $\|\bold X\|\leq h$, $\|\bold Y\|\leq h$


($h$ --- положительная постоянная), для каждого фиксированного значения


$k$ непрерывны как функции переменных $\bold X$, $\bold Y$ и


удовлетворяют условиям


$$


\bold G(k, \bold 0)= \bold 0,


\quad \|\bold R(k,\bold X,\bold Y)\|\leq


\psi(\bold X,\bold Y) \, \|\bold X\|^\mu,


\quad \|\bold D(k,\bold X,\bold Y)\|\leq c \|\bold X\|^\lambda,


$$


где $\psi(\bold X,\bold Y) \to 0$ при


$\|\bold X\|+\|\bold Y\|\to 0$, $c>0$, $\lambda>0$.





Наряду с уравнениями (26) рассмотрим систему, составленную из


уравнений (1) и уравнений


\begin{equation}


\bold Y(k+1)=\bold G(k,\bold Y(k)).


\end{equation}


Систему уравнений (1), (27) будем называть системой первого


приближения для~(26).





Исследуем вопрос, при каких условиях устойчивость нулевого решения


системы первого приближения обеспечивает устойчивость нулевого решения


системы~(26).





По-прежнему будем считать, что нулевое решение уравнений (1)


асимптотически устойчиво, причем для этих уравнений существуют


однородные функции Ляпунова $V(\bold X)$ и $W(\bold X)$,


обладающие свойствами 1)--3).


Кроме того, предположим, что для уравнений (27) найдена функция Ляпунова


$V_1(k,\bold Y)$, которая задана при $k=0,1,\dotsc$, $\|\bold Y\|\leq h$,


положительно определена, непрерывно дифференцируема по


$\bold Y$, ее частные производные


$\partial V_1/\partial y_j$, $j=1,\dotsc,l$, ограничены, а приращение


$V_1(k,\bold Y)$ на решениях уравнений (27) неположительно.


Заметим, что если функция $V_1(k,\bold Y)$ удовлетворяет указанным


условиям, то нулевое решение системы (27) является равномерно


устойчивым ([2], c.\,27).





\begin{thm}


При выполнении неравенства $\lambda> \mu-1$ нулевое решение 


системы $(26)$ устойчиво по всем переменным и асимптотически


устойчиво относительно $\bold X$.


\end{thm}





\begin{pf}


Рассмотрим приращения функций $V(\bold X)$


и $V_1(k, \bold Y)$ на решениях системы~(26). Получим, что при всех


$k=0,1,\dotsc$, $\|\bold X(k)\|\leq h$, $\|\bold Y(k)\|\leq h$


имеют место соотношения


\begin{gather*}


V(\bold X(k+1))- V(\bold X(k))\leq \|\bold X(k)\|^{m+\mu-1}


\big(-b_1+b_2 \, \psi(\bold X(k),\bold Y(k))\big), \\


V_1(k+1, \bold Y(k+1)) - V_1(k, \bold Y(k))\leq b_3 \,


\|\bold X(k)\|^\lambda,


\end{gather*}


где $b_1$, $b_2$, $b_3$ --- положительные постоянные.


Следовательно, существует такое число $\gamma> 0$, что если


при $k=k_0,\dotsc,k_1$ решение $(\bold X^\top(k), \bold Y^\top(k))^\top$


системы~(26) остается


в области $\|\bold X\|< \gamma$, $\|\bold Y\|< \gamma$,


то для этих значений $k$ справедлива оценка


$$


V(\bold X(k+1))\leq V(\bold X(k))-\frac{b_1}2 \|\bold X(k)\|^{m+\mu-1}.


$$





Учитывая неравенства (8), которым удовлетворяет однородная функция


$V(\bold X)$, имеем


$$


V(\bold X(k+1))\leq V(\bold X(k))-\frac{b_1}2


\bigg(\frac{V(\bold X(k))}{a_2}\bigg)^{\frac{\mu-1}m+1}.


$$


При этом число $\gamma$ выберем настолько малым,


чтобы выполнялось условие


$\frac{(\mu-1+m) b_1 \gamma^{\mu-1}}{2 m a_2}\leq 1$.





Используя лемму 2 и оценки (8), приходим к неравенству


\begin{equation}


\|\bold X(k)\| \leq c_1 \|\bold X(k_0)\| \big(1+c_2


\|\bold X(k_0)\|^{\mu-1} (k-k_0)\big)^{-\frac1{\mu-1}},


\end{equation}


которое справедливо при $k=k_0,\dotsc,k_1+1$.


Здесь $c_1$, $c_2$ --- положительные постоянные, не зависящие от


начальных данных решения.





Далее получаем


\begin{gather*}


V_1(k_1+1, \bold Y(k_1+1)) - V_1(k_0, \bold Y(k_0))\leq


b_3 \sum_{k=k_0}^{k_1} \|\bold X(k)\|^\lambda\leq \\


%


\leq b_3 \|\bold X(k_0)\|^\lambda \bigg(1+c^\lambda_1


\sum_{k=k_0+1}^{k_1} \big(1+c_2


\|\bold X(k_0)\|^{\mu-1}(k-k_0)\big)^{-\frac\lambda{\mu-1}}\bigg)\leq \\


%


\leq b_3 \|\bold X(k_0)\|^\lambda \bigg(1+


c^\lambda_1 \int_{k_0}^{k_1-1} \big(1+c_2


\|\bold X(k_0)\|^{\mu-1} (t-k_0)\big)^{-\frac\lambda{\mu-1}} dt \bigg).


\end{gather*}





Значит, имеет место оценка


\begin{equation}


V_1(k_1+1, \bold Y(k_1+1)) \leq V_1(k_0, \bold Y(k_0)) +


b_3 \|\bold X(k_0)\|^\lambda+b_3 c^\lambda_1 c_3 


\|\bold X(k_0)\|^{\lambda-\mu+1},


\end{equation}


где $c_3=\int\limits_{0}^{+\infty}(1+c_2 \tau


)^{-\frac\lambda{\mu-1}}d\tau$.





Зададим сколь угодно малое положительное число


$\varepsilon<\gamma$. Пусть


$\eta=\inf\limits_{k\geq 0, \, \|\bold Y\|=\varepsilon} V_1(k, \bold Y)$.


Выберем $\delta>0$ так, чтобы выполнялись условия


$c_1 \delta <\varepsilon$, $3b_3\delta^{\lambda} <\eta$,


$3b_3 c^\lambda_1 c_3\delta^{\lambda-\mu+1} <\eta$,


$V_1(k, \bold Y)<\eta/3$ при $\|\bold Y\|<\delta$, $k=0,1,\dotsc$





Из оценок (28) и (29) следует, что если начальные данные решения


$(\bold X^\top(k), \bold Y^\top(k))^\top$ удовлетворяют неравенствам


$k_0\geq 0$, $\|\bold X(k_0)\|< \delta$, $\|\bold Y(k_0)\|< \delta$,


то при всех $k\geq k_0$ имеем


$\|\bold X(k)\|< \varepsilon$, $\|\bold Y(k)\|< \varepsilon$, и при этом


$\|\bold X(k)\| \to 0$ при $k\to +\infty$.


\end{pf}





\begin{note}


Теорема 5 представляет собой распространение


теоремы Ляпунова--Малкина ([5], c.\,108--113) об асимптотической


устойчивости по отношению к части переменных на


случай разностных систем с существенно нелинейными уравнениями


первого приближения.


\end{note}





\begin{exam}


Пусть задано скалярное уравнение


\begin{equation}


\Ddot z+a(z) \Dot z{}^\sigma=u,


\end{equation}


где функция $a(z)$, $a(0)=0$, определена и непрерывна


при $|z|\leq h$ ($h$ --- положительная постоянная);


$\sigma$ --- рациональное число с нечетным знаменателем, $\sigma\geq 1$;


$u$ --- управление. Рассмотрим задачу стабилизации многообразия


$\Dot z=0$ уравнения (30) дискретным регулятором вида


$u(t)=\varphi(\Dot z(k))$ при $t\in [k, k+1)$, $k=0,1,\dotsc$





Исследуемое уравнение эквивалентно системе


\begin{equation}


\Dot x=-a(y) x^\sigma+u, \quad \Dot y=x.


\end{equation}


Интегрируя эту систему на промежутках $[k, k+1]$, приходим к системе


разностных уравнений


\begin{equation}


\begin{split}


x(k+1)&=x(k)+\varphi(x(k))+R(x(k),y(k)), \\


y(k+1)&=y(k)+D(x(k),y(k)),


\end{split}


\end{equation}


где $R(x(k),y(k))=-\int\limits_k^{k+1} a(y(t))x^\sigma(t) dt$,


$D(x(k),y(k))=\int\limits_k^{k+1} x(t) dt$.


Заметим, что если число $\delta$ достаточно мало, то при


$|x(k)|{<}\delta$, $|y(k)|{<}\delta$ справедливы оценки


$|R(x(k),y(k))|\leq \psi(x(k),y(k))|x(k)|^\sigma$,


$|D(x(k),y(k))|\leq c |x(k)|$.


Здесь $\psi(x,y) \to 0$ при $|x|+|y| \to 0$, $c>0$.





Пусть $\varphi(x(k))=-b x^\mu(k)$, где $b$ --- положительная постоянная,


$\mu>1$ --- рациональное число с нечетными числителем и знаменателем.


Применяя теорему~5, получаем, что при выполнении неравенств


$\mu<2$, $\sigma\geq \mu$ нулевое решение разностной системы (32)


устойчиво по всем переменным и асимптотически устойчиво относительно~$x$.


Но тогда тем же самым свойством обладает и нулевое решение системы


дифференциальных уравнений~(31).


\end{exam}








\begin{thebibliography}{9}








\bibitem{1}


Зубов В.И. {\it Проблема устойчивости процессов


управления.} -- Л.: Судостроение, 1980. -- 253~c.


\bibitem{2}


Халанай А., Векслер Д. {\it Качественная теория


импульсных систем.} -- М.: Мир, 1971. -- 310~c.


\bibitem{3}


Зубов В.И. {\it Устойчивость движения. Методы Ляпунова


и их применение.} Учебн. пособие. -- М.:


Высш. школа, 1973. -- 271~c.


\bibitem{4}


Бромберг~П.В. {\it Матричные методы в теории релейного


и импульсного регулирования.} -- М.: Наука, 1967. -- 324~c.


\bibitem{5}


Малкин И.Г. {\it Теория устойчивости движения.}


-- М.--Л.: Гостехиздат, 1952. -- 432~c.





\end{thebibliography}





\vskip 2cm





\post{\it Санкт-Петербургский}{\it Поступила}


\post{\it государственный университет}{25.04.2002}





\label{end}


\end{document}


